\section{Introduction}
\label{sec:intro}

\IEEEPARstart{L}{arge} language models (LLMs) are now widely deployed
as general-purpose generators, and the rise of tool use has composed
them into multi-step agent
loops~\cite{yao_react_2023,wang_survey_2024} that plan, invoke
external tools, and resume the same task after each action. Although
multi-turn chat already retains a conversation prefix in the
key-value (KV) cache, an agent loop further superimposes structured
tool waits, larger tool-result inserts, and more model invocations
per task, so that the object of serving is no longer a single request
but a session that is expected to return.

Serving systems have accordingly devoted most of their effort to how
tokens are computed, including how prefill and decode share or split
devices~\cite{patel_splitwise_2024,zhong_distserve_2024}. Under agent
loops, however, the end-to-end cost is often set by how the KV cache
is retained, evicted, and placed. Whereas a multi-turn conversation
remains paced by one user, agent workloads place many agents and many
sessions on the same pool at once and hold a growing prefix across
tool waits. The resulting live set shares static random-access memory (SRAM) and
high-bandwidth memory (HBM) with other attention clients, including
prefill, decode, and an optional vision-language model, on both edge
systems-on-chip (SoCs) and cloud servers. Once aggregate residency
outgrows that pool, eviction and hierarchical placement become an
efficiency problem in their own right and stand orthogonal to
compute-mode optimization. Fig.~\ref{fig:motivation}(a) records this
operating point.

\begin{figure}[t]
\centering
\includegraphics[width=\columnwidth]{fig01_motivation.pdf}
\caption{Motivation. \textbf{(a)}~Agent-resident occupancy versus a
request-scoped counterfactual on the same twelve-turn session.
\textbf{(b)}~Monotonic prefix growth across 1\,415 sessions into edge
and cloud SRAM/HBM envelopes. \textbf{(c)}~Recency, timeout, and ETA
remain above a B\'el\'ady oracle in AMAT, extra prefill, and TTFT.}
\label{fig:motivation}
\end{figure}

Fig.~\ref{fig:motivation}(a) contrasts two KV contracts on the same
twelve-turn session. The agent-resident curve holds the prefix through
every tool wait, whereas the request-scoped pulses drop the cache
after each decode. Unlike a multi-turn chat staircase whose idle
interval is a human pause that may never return, an agent tool wait is
a scheduled return, and recency or timeout policies therefore treat the
second kind of idle as if it were the first.

That resident contract becomes expensive once the prefix outgrows the
finite pool. Fig.~\ref{fig:motivation}(b) overlays prefix tokens on
coding and web-agent traces across 1\,415 sessions and three model
families. The curves rise monotonically into the SRAM and HBM
envelopes that edge and cloud stacks provision, and a budget that
cannot hold the live set is what makes session-level scheduling
unavoidable.

Published proxies miss the mechanism information that the loop
already produces and therefore treat a live wait as a cold,
discardable unit. Recency and timeout invert the rank by reading a
tool wait as coldness. A per-request lifecycle issues an absolute
keep-or-drop without ever ranking the live set against other
residents~\cite{abhyankar_infercept_2024,li_continuum_2026}. A
tool-type arrival table ranks residents yet assigns the same wait to
every job that shares a tool. Identity- or workflow-conditioned
scores separate residents only while roles or graphs remain diverse
and flatten toward recency once the batch is
uniform~\cite{cachescout_2026,pan_kvflow_2025,zheng_pbkv_2026}.
Hierarchical stores enlarge the pool yet still place from a static
plan, a job-scheduler hint, or prefix popularity, so that a kept
prefix is served on the slow path. Paged and radix
runtimes~\cite{kwon_pagedattention_2023,zheng_sglang_2024} manage the
pool at block granularity. Token, device, and interconnect engines
decide which bytes of one request to compute or move, a decision
that is complementary to ranking who stays in the shared pool, as
Table~\ref{tab:hw-kv} places on one grid. As Fig.~\ref{fig:motivation}(c)
reports, recency, timeout, and estimated time of arrival (ETA) all
sit above a B\'el\'ady next-reference oracle in average memory
access time (AMAT), extra prefill, and time to first token (TTFT)
across six traces and six capacity envelopes. Because that oracle is not observable online, a
useful substitute must exploit a signal that is already present on
every session, observable at runtime, and cheap to maintain, namely
the session's own gap history and its own progress toward completion.

Orthogonal to token datapaths and cluster routers that pursue the
same end-to-end efficiency from complementary vantage points, this
work proposes a near-memory scheduler beside the hierarchy
for long-horizon agent tasks in which many attention clients compete
for a two-tier pool on the edge and in the cloud. The scheduler
must decide who leaves the pool and who sits in the fast tier, and
both decisions are conditioned on the loop's mechanism information
rather than on an attempt to read task content.

To address those issues, we present
\textbf{U}nified \textbf{N}ative \textbf{I}nter-turn
\textbf{S}ession \textbf{O}rchestration \textbf{N}exus (\unison),
an event-driven near-memory scheduler in which
\textbf{S}urvival-\textbf{P}enalty \textbf{E}viction for
\textbf{A}gent \textbf{R}eturn-gap (\spear) and
\textbf{T}iering in \textbf{I}dle-window \textbf{D}MA
\textbf{E}vents (\tide) share one live ranking in a session
register file beside the memory hierarchy.
\spear scores remaining distance from a gap exponential moving
average (EMA) and a turn-indexed hazard and selects who leaves, while
\tide reads the complementary order, spends the observed wait as a
direct memory access (DMA) budget, and selects who sits in the appropriate memory hierarchy.
The main contributions are as follows.
\begin{enumerate}
\item \textbf{Algorithm.} \spear and \tide jointly rank sessions from
loop signals that the serving runtime already observes. Evaluated on
six traces that cross the
SWE-bench~\cite{jimenez_swebench_2024} and
GAIA~\cite{mialon_gaia_2023} agent benchmarks with three model
families~\cite{qwen3_2025,devstral_2025,gemma4_2026}, totaling 1\,415
sessions, 33\,596 turns, and six capacity envelopes from edge to
cloud, the joint policy is the best non-oracle entry on hit rate and
AMAT on every trace, raising the hit rate by $0.3\%$ to $23.1\%$ and reducing AMAT by
$22\%$ to $51\%$. On the long-horizon
traces it also lowers serving TTFT by 58\% to 89\% when extra
prefill sits on the critical path.
\item \textbf{Architecture.} We show that cycle-level event fidelity,
unified eviction and migration state, and deterministic gap-window
completion jointly require a near-memory control plane that cannot be
decomposed into independent IPs or realized in software. The 28-nm
CMOS scheduling core occupies 0.169\,mm$^{2}$ at 13.6\,mW and 150\,MHz for 64
sessions, a negligible overhead that reproduces the floating-point
ranking at Kendall $\tau{>}0.998$.
\end{enumerate}
